% =====================================================================
% Section: Introduction (with Related Work)  (outline.md §1)
% Draft v1 (2026-06-10). Related Work is integrated into the Introduction
% following the CS/engineering SCI body convention (outline.md).
% Citation keys refer to paper/references.bib.
% This draft contains NO quantitative result assertions; all comparative
% statements are qualitative (motivation / positioning) and are backed by
% the experiments referenced in plan/review/method-experiment-traceability.md.
% =====================================================================
\section{Introduction}
\label{sec:intro}

Single image super-resolution (SR) aims to reconstruct a high-resolution (HR)
image from a low-resolution (LR) observation, a problem that underpins applications
from mobile photography to display rendering. Before deep learning, SR was often
formulated with hand-crafted image priors, such as internal self-similarity
across scales~\cite{glasner2009single}. Since the seminal deep
model SRCNN~\cite{dong2014srcnn}, accuracy has improved through deeper,
residual, and feature-reuse networks, including
VDSR~\cite{kim2016vdsr}, RDN~\cite{zhang2018rdn}, RCAN~\cite{zhang2018rcan},
and EDSR~\cite{lim2017edsr}. A recent survey summarizes this progression across
architectures, losses, datasets, and metrics~\cite{moser2023survey}. The
corresponding growth in parameters and computation, however, makes many
high-capacity models impractical for mobile and edge devices, where memory,
latency, and energy are tightly constrained. This tension has motivated a sustained
line of \emph{lightweight} SR, which seeks a favorable
trade-off between reconstruction quality and model cost.

\paragraph{Lightweight CNN-based SR.}
Efficient CNN designs reduce cost by moving computation to low-resolution
features, reusing parameters, simplifying feature distillation, or replacing
heavy operators with compact mixing. CARN employs cascading
residual connections with group convolutions~\cite{ahn2018carn}; IMDN introduces
information multi-distillation~\cite{hui2019imdn}; RFDN reformulates this idea
with residual feature distillation blocks~\cite{liu2020rfdn}; and RLFN further
simplifies local feature extraction for efficient SR~\cite{kong2022rlfn}. More
recent lightweight designs also explore shuffle-based convolutional
mixing~\cite{sun2022shufflemixer} and dual-regression
constraints~\cite{guo2022dualregression}.
Efficient-SR challenges reflect the same deployment pressure by ranking models
with accuracy, runtime, parameters, and FLOPs considered together~\cite{ren2024ntire}.
For the present work, the main limitation of these methods is not efficiency but
context: locally connected convolutional operators cannot directly compare spatially
distant yet semantically similar regions without stacking many layers.

\paragraph{Lightweight Transformer-based SR.}
Long-range modeling tackles the context limitation from a different direction.
Non-local operators aggregate features across distant positions~\cite{wang2018nonlocal},
and Transformer-based restoration extends this idea with learned attention.
Swin Transformer provides the shifted-window
operator that later restoration models adopt for tractable attention~\cite{liu2021swin}.
SwinIR introduces window-based attention for image restoration~\cite{liang2021swinir};
ELAN streamlines long-range attention to cut redundant computation~\cite{zhang2022elan};
SRFormer revisits the window--channel trade-off with permuted
self-attention~\cite{zhou2023srformer}; and HAT activates more pixels to enlarge the
effective receptive field~\cite{chen2023hat}. Efficient variants further mix local
operators with omni-dimensional attention~\cite{wang2023omnisr} or n-gram
context~\cite{choi2023ngswin}, adopt mixed Transformer
blocks~\cite{zheng2023emt}, or integrate features hierarchically from coarse to
fine~\cite{liu2022hpinet}. More recent lightweight
Transformers push this efficiency further: GRFormer introduces grouped residual
self-attention to trim redundant parameters and computation~\cite{li2024grformer},
LMLT applies multi-level attention with per-head feature sizes to widen the
effective context~\cite{kim2025lmlt}, and Compacter shares a single projection
across spatial and channel attention for compact restoration~\cite{wu2024compacter}.
ATD takes a step toward content adaptivity by augmenting local windows with an
adaptive token dictionary~\cite{zhang2024atd}. Yet all of these still bind token
interaction to a predefined spatial partition: semantically similar but spatially
distant tokens interact only when the window grows or an extra global path is
added, which motivates organizing tokens by content instead of position.

\paragraph{Content-adaptive routing and token clustering.}
A complementary direction routes computation or interaction by \emph{content}
rather than by a fixed spatial layout. One thread allocates compute where
reconstruction is hard: PCSR classifies each pixel and dispatches it to an
upsampler of matching capacity~\cite{jeong2024pcsr}, while SalDRN treats visual
saliency as a proxy for region-level difficulty and selects feature paths of
appropriate depth accordingly~\cite{wu2023saldrn}. Pushing this idea to
real-world SR, MoR-DASR casts restoration as a mixture of experts whose
degradation-aware gate activates only relevant ranks and relies on a
load-balancing loss to stop expert usage from collapsing onto a few
branches~\cite{he2026mordasr}. A second thread instead groups \emph{tokens} by
content, so that distant but semantically related regions interact within a
compact attention---an idea rooted in non-local and self-exemplar SR that mines
recurring patches within an image~\cite{mei2020csnln} or transfers texture from a
reference~\cite{yang2020ttn}. SPIN clusters tokens into superpixels for
lightweight interaction~\cite{zhang2024spin}, and CATANet, the most direct
predecessor of our work, introduces a Token Aggregation Block (TAB) that keeps a
small set of routing centers, assigns each token to its nearest center, reorders
tokens by cluster label, and aggregates the reordered sequence with an
intra/inter-group attention operator~\cite{liu2025catanet}. This content routing
attains strong accuracy at a low budget; crucially, the load-balancing concern
raised for expert routing reappears here as a risk of \emph{prototype collapse}
that any learnable token router must also control.

\paragraph{Limitations of center-based routing.}
Despite its effectiveness, the TAB routing mechanism has four limitations.
(i) Its routing centers are maintained across batches through a $k$-means--style
iteration with an exponential moving average (EMA) buffer, so the centers are
driven by a cross-batch historical state rather than strictly by the current
image and adapt only slowly to per-image semantics. (ii) The token-to-center
relation is a hard \texttt{argmax} label that discards the \emph{confidence} of
the assignment, treating a token that clearly belongs to one cluster identically
to an ambiguous one on a boundary. (iii) The reordering clusters tokens only by
label, with no intra-cluster purity constraint, so high- and low-confidence tokens
are interleaved within a group during aggregation. (iv) Prototype \emph{generation}
and \emph{confirmation} are entangled in a single iterative update, which weakens
interpretability and stability. These observations suggest that a routing module
built directly from the current input, that preserves and exploits assignment
confidence, and that separates prototype generation from confirmation could
improve content routing without sacrificing the lightweight regime---provided it
also guards against the prototype collapse that a learnable router can induce, in
the same spirit as the load balancing used for expert routing~\cite{he2026mordasr}.

\paragraph{Our approach.}
We propose \emph{Dynamic Prototype Routing} (DPR), a drop-in replacement for the
TAB routing that keeps the CATANet backbone and its external tensor interface
unchanged while remaining in a comparable lightweight operating regime. DPR builds
semantic prototypes directly from the current tokens via a soft assignment. It then
confirms them with a learnable query bank, matches tokens back to the confirmed
prototypes with a learnable temperature, and propagates an explicit per-token
confidence through ordering and aggregation. The resulting network,
\textbf{DPRNet}, is trained on DIV2K and evaluated on five standard benchmarks at
$\times2/\times3/\times4$, where it maintains CATANet-level reconstruction quality
with small gains on most benchmark cells and a comparable lightweight cost profile,
while exposing an interpretable, confidence-aware routing
mechanism.

\paragraph{Contributions.}
The main contributions are as follows.
\begin{itemize}
  \item \textbf{C1 -- Dynamic semantic prototypes.} We generate routing prototypes
  from the current tokens by a soft, assignment-weighted aggregation, removing the
  cross-batch EMA center buffer so that prototypes follow per-image semantics
  (Section~\ref{sec:method_dpr}, Eqs.~\eqref{eq:assign}--\eqref{eq:proto_content}).
  \item \textbf{C2 -- Decoupled prototype confirmation.} We separate prototype
  \emph{generation} from \emph{confirmation} by refining the content prototypes
  with a learnable query bank through a single gated cross-attention, improving the
  stability and interpretability of the routing slots
  (Eq.~\eqref{eq:proto_refine}).
  \item \textbf{C3 -- Confidence-aware routing.} We explicitly retain the token
  assignment confidence and use it in three places: a composite ordering key that
  sorts high-confidence tokens first within each cluster, a confidence-modulated
  global gate, and a soft prototype fallback for uncertain tokens
  (Eqs.~\eqref{eq:sortkey}, \eqref{eq:iasa_gate}, \eqref{eq:soft_fallback}).
  \item \textbf{C4 -- Stabilized routing.} We introduce a learnable routing
  temperature together with a prototype usage-balancing regularizer to mitigate
  collapse onto a few prototypes and keep slot usage spread out
  (Eqs.~\eqref{eq:scores}, \eqref{eq:balance}).
\end{itemize}
Each design element is exposed as an independent switch that defaults to the full
model. Turning these switches off gives the corresponding ablated variants, which
allows the contributions to be isolated cleanly in the ablation study
(Section~\ref{sec:ablation}).
